% Responsable: Integrante B (enunciado, seccion 2.3)
% Ensayos, modelo con condiciones de validez, validacion con datos nuevos,
% polos y ceros, frecuencia de muestreo justificada.
\section{Identificaci\'on del modelo}

% Condiciones de cada escalón y resolución de la medición.
El protocolo contempla dos velocidades iniciales y tres niveles de
frenado normalizado $u\in[0,1]$ por velocidad. En cada ensayo a lazo
abierto se corta la alimentaci\'on del motor, se observa la
desaceleraci\'on libre y se aplica un escal\'on de frenado. El registro
conserva el tiempo, la orden $u$, la velocidad angular $\omega$ y
los eventos del ensayo. El periodo nominal del lazo es de
$20\,\mathrm{ms}$, igual al periodo del PWM del servo. La resoluci\'on
efectiva de la velocidad depende del tiempo entre marcas del disco.
El ajuste usa \'unicamente mediciones nuevas y compara la velocidad
registrada con el promedio predicho en cada intervalo entre marcas.

% La transferencia describe el aporte del freno respecto a la caída libre.
Para cada velocidad inicial se ajusta por m\'inimos cuadrados el modelo
incremental
\begin{equation}
  G(s)=\frac{\Delta\Omega(s)}{\Delta U(s)}
       =-\frac{K e^{-Ls}}{\tau s+1},
  \label{eq:planta-frenado}
\end{equation}
donde $K$ es la magnitud de la ganancia, $\tau$ es la constante de
tiempo y $L$ es el retardo efectivo. El signo negativo representa la
reducci\'on de velocidad al aumentar el frenado. El modelo tiene un
polo en $-1/\tau$ y carece de ceros finitos. Se estima una pendiente
com\'un de desaceleraci\'on libre para las capturas de cada velocidad
inicial y se ajusta una velocidad de partida por ensayo. Las ca\'idas
estimadas por nivel permiten buscar indicios de zona muerta y
saturaci\'on. El retardo estimado incluye el efecto de la resoluci\'on
temporal de la medici\'on.

% Las repeticiones reservadas se usan solo para evaluar las predicciones.
La validaci\'on conserva los par\'ametros identificados y fija la
velocidad de partida de cada repetici\'on con los datos previos a
aplicar el freno. Las predicciones se comparan con las mediciones
posteriores al escal\'on mediante el error absoluto medio (MAE) y la
ra\'iz del error cuadr\'atico medio (RMSE). La comparaci\'on se limita
al rango de frenado ensayado, con el motor apagado y el disco a\'un
en giro.

% Completar con parámetros, métricas y curvas al disponer de capturas reales.
La estimaci\'on experimental de los par\'ametros y su validaci\'on
quedan pendientes de las mediciones del banco.
